\chapter{Simulated Annealing and Its Parallel Variants}

The energy-based formulation of Chapter~1 reduces polygon packing to
minimizing a highly non-convex function $E(\conf; L)$ over
$\mathcal{C} = \mathbb{R}^{2N} \times [0, 2\pi)^N$.
This chapter develops the stochastic optimization machinery used
to explore that landscape.
We begin by establishing the common probabilistic target, then
derive the base Markov chain dynamics and the three parallel coordination
strategies---MS, ER-MS, and PT---as distinct approximations to sampling
from that target.

% ─────────────────────────────────────────────────────────────────
\section{The Common Probabilistic Target}
% ─────────────────────────────────────────────────────────────────

All three methods studied in this work can be understood as strategies for
approximating samples from the \emph{Gibbs measure} at temperature $T > 0$:
\[
  \pi_T(\conf) \;=\; \frac{1}{Z(T)}\exp\!\left(-\frac{E(\conf;L)}{T}\right),
  \qquad
  Z(T) = \int_{\mathcal{C}} e^{-E(\conf;L)/T}\,d\conf.
\]
As $T \to 0$, $\pi_T$ concentrates on the global minimizers of $E$.
Feasible configurations correspond to zero-energy states, so finding a
feasible packing is equivalent to sampling from $\pi_T$ in the low-temperature
limit.

The methods differ in how they allocate $R$ parallel workers to approximate
this target:
\begin{itemize}
  \item \textbf{MS} draws $R$ independent samples from separate realizations
        of the annealing process.
  \item \textbf{ER-MS} couples workers via a periodic selection--mutation
        operator, forming an interacting particle system that biases mass
        toward low-energy regions.
  \item \textbf{PT} couples $R$ chains thermodynamically via swap moves,
        targeting the \emph{joint} Gibbs distribution over the replica ensemble.
\end{itemize}

% ─────────────────────────────────────────────────────────────────
\section{Base Markov Chain Dynamics}
% ─────────────────────────────────────────────────────────────────

\subsection{Metropolis--Hastings Kernel}

Given a configuration $\conf_k$ at step $k$, a candidate $\conf'$ is drawn
from a symmetric Gaussian proposal kernel $q(\,\cdot \mid \conf_k)$ and
accepted with probability
\[
  \alpha(\conf_k, \conf')
  \;=\; \min\!\left(1,\; e^{-\Delta E / T_k}\right),
  \qquad
  \Delta E = E(\conf'; L) - E(\conf_k; L).
\]
Symmetry of $q$ ensures this satisfies detailed balance with respect to
$\pi_{T_k}$. The resulting transition kernel is
\[
  P_k(\conf \to d\conf')
  = q(\conf' \mid \conf)\,\alpha(\conf, \conf')\,d\conf'
  + r(\conf)\,\mathbf{1}_{\conf}(d\conf'),
\]
where $r(\conf) = 1 - \int q(\conf'\mid\conf)\,\alpha(\conf,\conf')\,d\conf'$
is the rejection probability.

\paragraph{Proposal structure.}
At each step, a polygon index $i \sim \mathrm{Uniform}\{1,\dots,N\}$ is
selected and a joint positional and rotational perturbation is applied:
\[
  x_i' = x_i + \delta_x, \quad
  y_i' = y_i + \delta_y, \quad
  \theta_i' = \theta_i + \delta_\theta,
\]
with $(\delta_x, \delta_y)^\top \sim \mathcal{N}(\mathbf{0},
\sigma_{\mathrm{pos}}^2 I_2)$ and
$\delta_\theta \sim \mathcal{N}(0, \sigma_\theta^2)$, drawn independently.

\subsection{Cooling Schedule and Mixing}

We use a \emph{geometric} cooling schedule:
\[
  T_k = T_{\max}\,\rho^k,
  \qquad
  \rho = \left(\frac{T_{\min}}{T_{\max}}\right)^{1/K},
\]
where $K$ is the total number of steps in a bisection probe.
The geometric schedule is characterized by a constant multiplicative
temperature reduction $\rho$ per step, which gives equal thermal exposure
at each decade of energy: the number of steps spent in any interval
$[cT, T]$ is $-\log c / \log(1/\rho)$, independent of the absolute
temperature level.

The theoretically sufficient condition for convergence to the global minimum
requires logarithmic cooling $T_k = c / \log(1 + k)$ for a constant $c$
exceeding the depth of the deepest local energy barrier~\cite{Hajek1988}.
Such schedules are computationally infeasible for the problem sizes considered
here; the geometric schedule instead prioritizes practical convergence within
a fixed computational budget, with $T_{\min}$ and $T_{\max}$ fixed by
pilot calibration (Section~\ref{sec:pt-pilot}).

\subsection{Integration with Bisection}

Each method operates within an outer bisection loop over the container
side length $L$. For a fixed bracket $[L_{\mathrm{lo}}, L_{\mathrm{hi}}]$,
a probe evaluates the midpoint $L = (L_{\mathrm{lo}} + L_{\mathrm{hi}})/2$
by running $K$ steps of the SA-based search.
If a feasible configuration is found, the bracket contracts:
$L_{\mathrm{hi}} \leftarrow L$; otherwise $L_{\mathrm{lo}} \leftarrow L$.
The SA energy function $E(\conf; L)$ is therefore re-evaluated at a
different $L$ on each bisection step, and the Gibbs target $\pi_T$
changes accordingly.

% ─────────────────────────────────────────────────────────────────
\section{Method A: Multi-Start Simulated Annealing (MS)}
\label{sec:ms}
% ─────────────────────────────────────────────────────────────────

\subsection{Definition}

MS instantiates $R$ independent Markov chains, each evolving under its
own copy of the base SA kernel:
\[
  \conf_k^{(r)} \;\xrightarrow{P_k,\, T_k}\; \conf_{k+1}^{(r)},
  \qquad r = 1, \dots, R.
\]
Workers share no state during a probe. Seeds are generated deterministically as
\[
  s_r = \mathrm{splitmix64}\!\bigl(
    g \oplus \mathrm{hash}(N, \texttt{run\_id}, \texttt{probe\_idx}, r)
  \bigr),
\]
where $g$ is the global seed, ensuring exact reproducibility.

\subsection{Probabilistic Analysis}

Let $p$ denote the probability that a single SA chain finds a feasible
configuration within the allotted probe time $\tau$.
The probability that at least one of $R$ independent workers succeeds is
\[
  P_{\mathrm{success}} = 1 - (1-p)^R.
\]
For small $p$, this grows as $Rp$ to first order, so $R$ workers provide
a linear improvement in success probability at linear cost.
More precisely, if the first-success time for a single chain is
$\mathrm{Geometric}(p)$-distributed with mean $1/p$, then the minimum
across $R$ independent chains has mean $1/(Rp)$: the expected time to
first feasibility scales as $1/R$.

\subsection{Limitations}

The independence of MS chains is both its computational strength
(zero communication overhead) and its fundamental weakness.
Workers cannot exploit configurations discovered by their siblings,
leading to redundant exploration of the same basins.
In highly multimodal landscapes with narrow feasible regions, this
redundancy becomes the dominant cost.

% ─────────────────────────────────────────────────────────────────
\section{Method B: Elite-Resample Multi-Start (ER-MS)}
\label{sec:erms}
% ─────────────────────────────────────────────────────────────────

\subsection{Motivation and Interpretation}

ER-MS augments MS by introducing a periodic selection--mutation operator
that concentrates the worker population toward low-energy regions.
Viewed abstractly, the $R$ workers form an \emph{interacting particle system}
whose empirical distribution
\[
  \mu_k^R = \frac{1}{R}\sum_{r=1}^R \delta_{\conf_k^{(r)}}
\]
is driven toward $\pi_T$ by alternating between independent SA evolution
and population-level reweighting. This mirrors the structure of
Sequential Monte Carlo (SMC) samplers, where a population of weighted
particles is evolved and periodically resampled.

\subsection{Resampling Operator}

A resampling event is triggered every
$K_{\mathrm{resample}} = 200N$ MH steps.
Workers are ranked by energy and the elite pool is defined as the top
\[
  k' = \lceil 0.25 R \rceil
\]
workers. Elite workers are preserved unchanged. Each non-elite worker $r$
is replaced by a perturbed copy of a uniformly drawn elite $r^*$:
\[
  \conf^{(r)} \;\leftarrow\; \conf^{(r^*)} + \bm{\xi}_r,
  \quad
  r^* \sim \mathrm{Uniform}\{1, \dots, k'\}.
\]

\subsection{Perturbation Model}

The perturbation $\bm{\xi}_r$ is drawn independently for spatial and
angular components.
For polygon $i$, the displacement and rotation noise are
\[
  (\delta x_i, \delta y_i)^\top \sim \mathcal{N}\!\left(\mathbf{0},\,
    \sigma_{\mathrm{pos}}(f)^2\, I_2\right),
  \qquad
  \delta\theta_i \sim \mathcal{N}\!\left(0,\, \sigma_\theta(f)^2\right),
\]
where the standard deviations decay linearly with the elapsed fraction
$f \in [0,1]$ of the current probe:
\begin{align}
  \sigma_{\mathrm{pos}}(f) &= (0.02 - 0.015\,f)\,L, \\
  \sigma_\theta(f) &= 5^\circ - 4^\circ f.
\end{align}
The full per-polygon perturbation covariance is therefore the block-diagonal matrix
\[
  \Sigma(f) = \mathrm{diag}\!\bigl(
    \underbrace{\sigma_{\mathrm{pos}}^2, \sigma_{\mathrm{pos}}^2,
    \sigma_\theta^2}_{\text{polygon }1},\,
    \dots,\,
    \underbrace{\sigma_{\mathrm{pos}}^2, \sigma_{\mathrm{pos}}^2,
    \sigma_\theta^2}_{\text{polygon }N}
  \bigr).
\]
Large early-probe noise allows workers to escape elite basins
and explore new regions; decaying late-probe noise concentrates
refinement near the current best solution.

\subsection{Anti-Collapse Mechanism}

Population collapse---all workers converging to the same local minimum---
reduces ER-MS to a single chain and eliminates the benefit of
parallelism. Collapse is detected via the normalized energy spread
\[
  \delta_{\mathrm{rel}}
  = \frac{E_{\max} - E_{\min}}{\max\!\left(1,\, \lvert E_{\mathrm{mean}}\rvert\right)}
  < 10^{-6},
\]
evaluated at each resampling event.
If this condition holds for two consecutive resamples, the perturbation
scales are doubled for the immediately following resample cycle:
\[
  \sigma_{\mathrm{pos}} \leftarrow 2\,\sigma_{\mathrm{pos}}(f), \qquad
  \sigma_\theta \leftarrow 2\,\sigma_\theta(f).
\]

\subsection{Relationship to SMC}

In standard SMC, particles carry importance weights and resampling is
performed proportional to those weights.
ER-MS replaces importance weights with a hard rank cutoff (top $k'$),
which is computationally cheaper but introduces selection bias:
the retained configurations are not guaranteed to be representative
samples from $\pi_T$.
The mutation step (Gaussian perturbation) plays the role of the MCMC
rejuvenation kernel in SMC, restoring diversity lost to selection.

% ─────────────────────────────────────────────────────────────────
\section{Method C: Parallel Tempering (PT)}
\label{sec:pt}
% ─────────────────────────────────────────────────────────────────

\subsection{Replica Ensemble and Joint Distribution}

PT maintains $R$ replicas at a fixed geometric temperature ladder:
\[
  T_r = T_{\min}\!\left(\frac{T_{\max}}{T_{\min}}\right)^{(r-1)/(R-1)},
  \qquad r = 1, \dots, R,
\]
so that $T_1 < T_2 < \cdots < T_R$.
The extended system targets the \emph{joint} Gibbs distribution
\[
  \Pi\!\left(\conf^{(1)}, \dots, \conf^{(R)}\right)
  = \prod_{r=1}^R \pi_{T_r}\!\left(\conf^{(r)}\right).
\]
The geometric ladder spacing gives a constant temperature ratio between
adjacent rungs, which empirically yields uniform swap acceptance rates
across the ladder.

\subsection{Within-Replica Dynamics}

Between swap events, each replica evolves independently under its own
MH kernel targeting $\pi_{T_r}$:
\[
  \conf_k^{(r)} \;\xrightarrow{P_k^{(r)},\, T_r}\; \conf_{k+1}^{(r)}.
\]

\subsection{Replica Exchange}

Every $K_{\mathrm{swap}} = 200N$ MH steps, an exchange is proposed
between each adjacent pair $(r, r+1)$.
The swap acceptance probability is
\[
  \alpha_{\mathrm{swap}}^{(r,r+1)}
  = \min\!\left(1,\;
    \exp\!\left[
      \left(\frac{1}{T_r} - \frac{1}{T_{r+1}}\right)
      \Bigl(E(\conf^{(r+1)};L) - E(\conf^{(r)};L)\Bigr)
    \right]
  \right).
\]

\begin{proposition}[Detailed Balance for Replica Exchange]
The swap acceptance probability $\alpha_{\mathrm{swap}}^{(r,r+1)}$
satisfies detailed balance with respect to the joint distribution $\Pi$.
\end{proposition}

\begin{proof}
Let $\bm{\conf}$ and $\bm{\conf}'$ be the ensemble states before and after
swapping replicas $r$ and $r+1$. The ratio of joint probabilities is
\begin{align*}
  \frac{\Pi(\bm{\conf}')}{\Pi(\bm{\conf})}
  &= \frac{\pi_{T_r}(\conf^{(r+1)})\,\pi_{T_{r+1}}(\conf^{(r)})}
         {\pi_{T_r}(\conf^{(r)})\,\pi_{T_{r+1}}(\conf^{(r+1)})}
   = \exp\!\left[
     \left(\frac{1}{T_r} - \frac{1}{T_{r+1}}\right)
     \left(E^{(r+1)} - E^{(r)}\right)
   \right],
\end{align*}
where $E^{(r)} = E(\conf^{(r)};L)$.
The Metropolis acceptance ratio $\alpha_{\mathrm{swap}}^{(r,r+1)} =
\min(1, \Pi(\bm{\conf}')/\Pi(\bm{\conf}))$ then satisfies
$\Pi(\bm{\conf})\,\alpha_{\mathrm{swap}}^{(r,r+1)}
= \Pi(\bm{\conf}')\,\alpha_{\mathrm{swap}}^{(r+1,r)}$
by the standard Metropolis argument.
\end{proof}

\subsection{Mechanism: Temperature-Assisted Basin Escape}

High-temperature replicas flatten the effective energy landscape,
allowing them to traverse barriers that would trap a cold chain.
When a swap carries a favorable configuration from replica $r+1$ to
replica $r$, that configuration begins cold-chain refinement from a
point already near a low-energy basin.
This \emph{temperature-assisted diffusion} is the core reason PT
outperforms MS on landscapes with high, narrow barriers separating
deep basins.

\subsection{Early Stopping Criterion}

Because our goal is feasibility rather than equilibrium sampling,
a probe terminates only when the \emph{coldest} replica ($r = 1$)
reaches $E(\conf^{(1)};L) \le \varepsilon_{\mathrm{feas}}$.
Feasible configurations visiting higher-temperature replicas are
not accepted as solutions unless they propagate to replica 1 via
an accepted swap chain.

\subsection{Pilot Calibration}
\label{sec:pt-pilot}

The temperature bounds $T_{\min}$ and $T_{\max}$ are set by the
following pilot procedure:
\begin{enumerate}
  \item Run a 10-minute pilot at $N=100$, $R=10$ with an initial
        temperature ladder.
  \item Compute the empirical swap acceptance rate $\hat{\alpha}_r$ at
        each adjacent pair $(r, r+1)$.
  \item If $\hat{\alpha}_r \notin [0.15,\, 0.60]$ for any rung,
        multiply $T_{\max}$ or $T_{\min}$ by a factor of 2 and repeat.
  \item Lock the calibrated values. Run a 5-minute sanity pilot at
        $N=200$ to confirm validity at larger system size.
\end{enumerate}
The target window $[0.15, 0.60]$ is a standard empirical criterion:
below 0.15 the cold chain mixes too slowly; above 0.60 adjacent
replicas are too similar for exchange to provide useful information flow.

% ─────────────────────────────────────────────────────────────────
\section{Energy Landscape Interpretation}
% ─────────────────────────────────────────────────────────────────

The relative performance of the three methods is determined by the
structure of the energy landscape $E(\cdot; L)$.
Three geometric features are particularly relevant:

\begin{itemize}
  \item \textbf{Local minima.} Configurations $\conf^*$ satisfying
        $\nabla E = 0$ with $E(\conf^*) > 0$ are jammed states: partially
        overlapping packings with no locally improving move.
  \item \textbf{Energy barriers.} A barrier of height $\Delta$ between
        two basins requires the chain to traverse a region of energy at
        least $\Delta$ above the current state. Under the Metropolis rule,
        the probability of acceptance at temperature $T$ scales as
        $e^{-\Delta/T}$, so barriers suppress transitions exponentially
        as $T \to 0$.
  \item \textbf{Basins of attraction.} Each local minimum is surrounded
        by a basin of configurations from which gradient-following
        dynamics converge to it. Wide, shallow basins are easily escaped
        at moderate temperatures; deep, narrow feasible basins require
        low temperatures to be refined once entered.
\end{itemize}

MS relies on independent discovery of distinct basins; it does not
improve mixing within any single chain.
ER-MS concentrates the worker population in the lowest-energy basins
discovered so far, amplifying exploitation at the cost of diversity.
PT enables direct inter-basin transitions by lifting configurations to
temperatures where barriers are effectively invisible, then returning
them to low temperature for refinement.

% ─────────────────────────────────────────────────────────────────
\section{Summary}
% ─────────────────────────────────────────────────────────────────

Table~\ref{tab:methods} summarizes the key structural properties of
the three strategies.

\begin{table}[h]
\centering
\caption{Structural comparison of parallel SA variants.}
\label{tab:methods}
\begin{tabular}{llll}
\toprule
Property & MS & ER-MS & PT \\
\midrule
Probabilistic target    & $\pi_T$ (each chain)        & Biased toward low $E$           & $\prod_r \pi_{T_r}$ \\
Inter-worker comm.      & None                        & Periodic (resample)             & Periodic (swap)     \\
Comm.\ trigger          & ---                         & Every $200N$ steps              & Every $200N$ steps  \\
Temperature structure   & Shared schedule             & Shared schedule                 & Fixed ladder        \\
Elite/cold bias         & None                        & Top $\lceil 0.25R\rceil$        & Cold replica only   \\
Barrier crossing        & Weak (single-chain)         & Moderate (population spread)    & Strong (PT diffusion) \\
Early stop condition    & Any worker feasible         & Any worker feasible             & Cold replica feasible \\
HPC pattern             & Embarrassingly parallel     & Barrier + collective sort       & Neighbor communication \\
\bottomrule
\end{tabular}
\end{table}

The three methods represent a progression from zero coupling (MS)
through unidirectional information flow driven by energy ranking (ER-MS)
to bidirectional thermodynamic coupling governed by detailed balance (PT).
Chapter~3 defines the energy function $E(\conf; L)$ explicitly and
develops the geometric algorithms required to evaluate it.